\documentclass[11pt]{article}
\usepackage[margin=1in]{geometry}
\usepackage[T1]{fontenc}
\usepackage{lmodern}
\usepackage{microtype}
\usepackage{xcolor}
\usepackage{hyperref}
\usepackage{booktabs}
\usepackage{enumitem}
\usepackage{listings}
\usepackage{tcolorbox}
\usepackage{parskip}
\usepackage{array}
\newcolumntype{L}[1]{>{\raggedright\arraybackslash}p{#1}}

% \path{} (from url, loaded by hyperref) prints paths and identifiers
% verbatim -- so `--` is not turned into a dash -- and may break at / _ - .
\makeatletter
\g@addto@macro\UrlBreaks{\do\_\do\-\do\.}
\makeatother
\urlstyle{tt}
\setlength{\emergencystretch}{3em}
\newcommand{\dd}{-{}-}

\hypersetup{
  colorlinks=true,
  linkcolor=blue!60!black,
  urlcolor=blue!60!black
}

\definecolor{codebg}{RGB}{245,245,245}
\definecolor{codecomment}{RGB}{90,100,90}
\definecolor{codekeyword}{RGB}{0,70,140}
\definecolor{codestring}{RGB}{150,60,20}

\lstdefinestyle{pythonstyle}{
  language=Python,
  backgroundcolor=\color{codebg},
  basicstyle=\ttfamily\small,
  keywordstyle=\color{codekeyword}\bfseries,
  commentstyle=\color{codecomment}\itshape,
  stringstyle=\color{codestring},
  showstringspaces=false,
  breaklines=true,
  frame=single,
  rulecolor=\color{black!20},
  columns=fullflexible,
  keepspaces=true
}
\lstdefinestyle{shellstyle}{
  language=bash,
  backgroundcolor=\color{codebg},
  basicstyle=\ttfamily\small,
  showstringspaces=false,
  breaklines=true,
  frame=single,
  rulecolor=\color{black!20},
  columns=fullflexible,
  keepspaces=true
}
\lstset{style=pythonstyle}

\newtcolorbox{keyidea}{
  colback=blue!5, colframe=blue!50!black,
  title=Key Idea, fonttitle=\bfseries
}
\newtcolorbox{tipbox}{
  colback=green!5, colframe=green!45!black,
  title=Tip, fonttitle=\bfseries
}
\newtcolorbox{warningbox}{
  colback=orange!7, colframe=orange!60!black,
  title=Common Pitfall, fonttitle=\bfseries
}
\title{
  \textbf{Codebase Change Summary --- Week 13}\\[4pt]
  \large Continuous Testing\\[2pt]
  \normalsize SE 413/513 --- Software Testing
}
\author{Demo Testing Suite}
\date{}

\begin{document}
\maketitle

\begin{center}
\textit{What changed in \path{demo-testing-suite/} (and the repository) since Week 11, and why}
\end{center}

\vspace{0.3cm}

% ============================================================
\section{Purpose of This Document}
% ============================================================
The \path{demo-testing-suite/} codebase is a single, shared project that
grows one lecture at a time. This document summarizes what was added or
changed for Week 13 and explains \emph{why} each change exists --- which
concept from the Week 13 reading, \textit{Continuous Testing: Every
Change, Every Time}, it demonstrates. Read it alongside the reading
(\path{reading_materials/w13_continuous_testing.pdf}) and
\path{demos/continuous_advanced_testing/README.md}.

\begin{keyidea}
Nothing from Weeks 1--11 was removed or changed in behavior. This week
adds one self-contained project, shared with Week 14: Week 11's shop
plus a \emph{pipeline} that runs every test level in order, cheapest
first, on a laptop and --- through a GitHub Actions workflow --- on
every push. The pipeline's first run found a real, latent defect.
\end{keyidea}

% ============================================================
\section{What Changed At a Glance}
% ============================================================
\begin{center}
\small
\begin{tabular}{@{}L{6.0cm}L{1.5cm}L{6.9cm}@{}}
\toprule
\textbf{File} & \textbf{Status} & \textbf{Purpose} \\
\midrule
\multicolumn{3}{@{}l}{\textit{The reading's project (not in the default run)}}\\
\path{demos/continuous_advanced_testing/} & New (copy) & Week 11's shop
and its 60 tests, copied; Week 11's folder is unchanged. \\
\path{pipeline.py} & New & The pipeline: \texttt{static},
\texttt{unit}, \texttt{integration} (80\% coverage gate),
\texttt{contract}, \texttt{system}, \texttt{acceptance}; stops at the
first failure; writes JUnit reports and keeps server logs. \\
\path{.coveragerc} & New & Coverage of \path{shop/} only, without the
defective variants; also measures the server processes. \\
\path{shop/payments.py} & Fixed & One line: pylint's E0606, found by the
pipeline's first run (Section~\ref{sec:finding}). \\
\path{flaky_examples/} & New & Two tests that are flaky on purpose: a
timing assumption and an order dependency. \\
\path{tools/rerun.py} & New & Runs the same tests N times and counts
the outcomes. \\
\midrule
\multicolumn{3}{@{}l}{\textit{Repository, tooling, and documentation}}\\
\path{.github/workflows/shop-pipeline.yml} & New & The pipeline in
GitHub Actions: a fast job on Python 3.12 and 3.13, then a slow job;
artifacts; a \emph{Run workflow} button with a choice of defect. \\
\path{.gitignore} & Modified & Ignores the pipeline's \path{reports/}. \\
\texttt{pytest.ini} & Modified & Registers the \texttt{contract} marker
(Week 14). \\
\path{docs/DEMO_GUIDE_week13_continuous_testing.md} & New &
Section-by-section lecture companion. \\
\texttt{README.md} (suite and repository) & Modified & Layout and the
Week 13 entry. \\
\bottomrule
\end{tabular}
\end{center}

\begin{tipbox}
No new dependencies: \texttt{pylint}, \texttt{pytest-cov}, and
\texttt{pytest-xdist} were added in Weeks 5--7. If your environment is
older than that, reinstall:
\begin{lstlisting}[style=shellstyle]
cd demo-testing-suite
pip install -r requirements.txt
\end{lstlisting}
\end{tipbox}

% ============================================================
\section{The Pipeline, Runnable}
% ============================================================
\textbf{Location:} \path{demos/continuous_advanced_testing/}. \textbf{Run
its commands from inside the folder:}
\begin{lstlisting}[style=shellstyle]
cd demos/continuous_advanced_testing
pytest                         # 67 passed (Week 11's 60, plus 7 contract tests)
python pipeline.py --list      # the six stages
python pipeline.py             # all six pass, about 12 seconds
\end{lstlisting}

Each defective variant stops the pipeline at the first stage that can
see it:

\begin{center}
\small
\begin{tabular}{@{}lll@{}}
\toprule
\textbf{Variant} & \textbf{Stops at} & \textbf{Result there}\\
\midrule
\texttt{SHOP\_SERVICE=racy} & system & 22/23 passed\\
\texttt{SHOP\_COUPONS=per\_unit} & acceptance & 14/16 passed\\
\texttt{SHOP\_PAYMENTS=no\_idempotency} (Week 14) & contract & 4/7 passed\\
\bottomrule
\end{tabular}
\end{center}

In GitHub the same stages run on every push and pull request that
touches the folder (about 50 seconds); once the workflow is on the
default branch, \emph{Actions $\rightarrow$ Shop pipeline $\rightarrow$
Run workflow} can run it against one of the defects.

% ============================================================
\section{A Real Finding: Use Before Definition}
\label{sec:finding}
% ============================================================
The pipeline's first stage runs \texttt{pylint \dd errors-only}. On its
first run it stopped at:
\begin{lstlisting}[style=shellstyle]
shop/payments.py:43:37: E0606: Possibly using variable 'problem' before assignment
\end{lstlisting}
With \texttt{retries=-1}, the payment client's \texttt{charge()} raised
\texttt{UnboundLocalError} instead of \texttt{PaymentUnavailable}: Week
6's use-before-definition anomaly, in the shop since Week 10, missed by
every test. It is fixed (one line) in the new folder only; Weeks 10 and
11's folders are left as their readings printed them.

% ============================================================
\section{Updated: Expected Test Results}
% ============================================================
From the project root:
\begin{lstlisting}[style=shellstyle]
cd demo-testing-suite
pytest -v
\end{lstlisting}
you should still see \textbf{75 passed, 3 xfailed, 1 failed}, exactly
Week 11's result. The new folder is not part of that run
(\texttt{pytest.ini} sets \texttt{testpaths = tests}), and the workflow
runs only the new folder: the main suite's deliberately failing test
would otherwise make every run red.

% ============================================================
\section{How This Connects to the Reading}
% ============================================================
\begin{keyidea}
Weeks 2--11 built tests at every level. Week 13 makes them run on every
change, in order of cost, on a clean machine, with gates and evidence
--- because a test that is not run finds nothing.
\end{keyidea}

% ============================================================
\section{Suggested Next Steps}
% ============================================================
\begin{enumerate}
\item Run \texttt{pipeline.py} against each defective variant, and
\emph{predict} the failing stage first.
\item Open the workflow file and a run in the \emph{Actions} tab side by
side.
\item Fix both flaky examples (Exercises 3 and 4), and prove it with
\path{tools/rerun.py}.
\item Add a pipeline to your semester project, and make it fail on
purpose at every stage.
\end{enumerate}

\end{document}
